\section*{A negative shared-data curvature refund at a vanishing selector gap}

All logarithms are base two.  Let $f_r$ be the OR of $k=2^r$ disjoint
$r$-bit point defects, and expose one bit $s$ of one block.  Put
$\epsilon=2^{-r}$, $P=(1-\epsilon)^{1/\epsilon}$, and
\[
 \delta=\operatorname{Inf}_s(f_r)
 =2\epsilon(1-\epsilon)^{k-1}={2P\epsilon\over1-\epsilon}. \tag{1}
\]
Write
\[
 Q_s=G(f_r)-{G((f_r)_{s,+})+G((f_r)_{s,-})\over2}.
\]
Commutation of restriction operators gives the exact complete increment
\[
 Q_s=\Delta_sU+\sum_{i\ne s}\Delta_s\Delta_iU.          \tag{2}
\]

\paragraph{Theorem.}
As $r\to\infty$,
\[
 {\Delta_sU\over\delta}=1+2r2^{-r}+O(2^{-r}),          \tag{3}
\]
\[
 {Q_s\over\delta}=1-2r^2 2^{-r}+O(r2^{-r}),            \tag{4}
\]
and therefore
\[
 {1\over\delta}\sum_{i\ne s}\Delta_s\Delta_iU
 =-2r^2 2^{-r}+O(r2^{-r}).                              \tag{5}
\]
In particular, for all sufficiently large $r$,
\[
 \Delta_sU>\operatorname{Inf}_s(f_r),\qquad
 Q_s<\operatorname{Inf}_s(f_r).                         \tag{6}
\]
Thus the selector-only payment is false on this vanishing-gap family, but
the full signed GU increment is paid by the same physical influence after
the shared-data curvature refund is retained.

\paragraph{Derivation.}
The exact heterogeneous upper-set formula from the preceding source is
\[
 U(\mathbf r)=c((2P-1)^2)+4P\sum_j\epsilon_j
 \left(r_j-{r_j\over2(1-\epsilon_j)}\right)
 +4P\sum_jL_j((1-\epsilon_j)-P),                        \tag{7}
\]
where $c(x)=x\log_2(1/x)$ and
$L_j=-\log_2(1-\epsilon_j)$.  The two physical children replace the
chosen $r$-block respectively by an $(r-1)$-block or delete it.  Hence
\[
 \Delta_sU=U_{k,r}-{U_{(r-1,r,\ldots,r)}+U_{k-1,r}\over2}. \tag{8}
\]
By block symmetry, $G(f_r)=kr\Delta_sU$, while the exact influence formula
is $I(f_r)=kr\delta$.  The refined growing-rank calculation gives
\[
 {G(f_r)\over I(f_r)}=1+2r\epsilon+O(\epsilon),         \tag{9}
\]
which is exactly (3), not an analogy.

To compute the mixed term without discarding the unequal children, set
$E=G-I$.  The influence recursion and (2) imply the exact identity
\[
 Q_s-\delta=E(f_r)-{E((f_r)_{s,+})+E((f_r)_{s,-})\over2}. \tag{10}
\]
Taylor expansion of (7) and its physical restrictions, uniformly for
$P$ in a fixed neighbourhood of $e^{-1}$, gives for a heterogeneous
OR-of-defects vector
\[
 E(\mathbf r)=4P\sum_j r_j^2\epsilon_j^2+mathcal E(\mathbf r), \tag{11}
\]
where changing one $r$-block to an $(r-1)$-block or deleting it changes the
remainder by $O(r\epsilon^2)$.  This follows directly from
$L_j=\epsilon_j/\ln2+O(\epsilon_j^2)$ and the bounded third derivatives of
$c((2P-1)^2)$ near $P=e^{-1}\ne1/2$; every first-order product change
cancels because
\[
 P_-={P\over1-\epsilon},\quad
 P_+={P(1-2\epsilon)\over1-\epsilon},\quad
 {P_-+P_+\over2}=P.                                     \tag{12}
\]

For the leading functional in (11), let
$S=kr^2\epsilon^2$.  Substitution of the two children is elementary and
gives
\[
 4PS-2(P_-S_-+P_+S_+)
 =-4Pr^2\epsilon^2+O(Pr\epsilon^2).                     \tag{13}
\]
Equations (10)--(13), divided by
$\delta=2P\epsilon/(1-\epsilon)$, prove (4).  Subtracting (3) from (4)
proves (5).

\paragraph{GU and terminal implication.}
For the frozen pair $(a,b)=(1,0)$ the local residual increment is
$Q_s-\delta<0$ for large $r$: this selector step has a genuine nonreused
reserve.  Nevertheless the companion B16 calculation proves
$G(f_r)-I(f_r)>0$.  Therefore the negative drift does not close GU unless
the surviving child/terminal residual is also paid.  This supplies an
explicit warning against dropping the $U$ boundary: favorable local drift
and a vanishing selector gap can coexist with a positive root residual.
It is a proved restricted mechanism, not a universal GU theorem.

